\newpage % Pour commencer la bibliographie sur une nouvelle page

% Ajoute une entrée "Bibliographie" dans la table des matières
\addcontentsline{toc}{chapter}{Bibliographie} 

\begin{thebibliography}{La Retraite en clair, s.d.}

\bibitem[CFDT Retraités, s.d.]{cfdt}
CFDT RETRAITÉS. « Histoire, origine et évolution des régimes de retraite ». \textit{Guide des pensions de retraite}, fiche 2. Consulté le 3 novembre 2026.

\bibitem[CNAV, 2015]{cnav2015}
CAISSE NATIONALE D'ASSURANCE VIEILLESSE. \textit{Abrégé statistique 2015}. Direction statistiques, prospective et recherche, 2015.

\bibitem[Conseil d'État, s.d.]{conseiletat}
CONSEIL D'ÉTAT. « Pierre Laroque ». \textit{Les grandes figures de la juridiction administrative}. Consulté le 3 novembre 2026.

\bibitem[COR, 2022]{cor2022}
CONSEIL D'ORIENTATION DES RETRAITES. « Histoire de la retraite : repères chronologiques ». Document n° 3, séance plénière du 15 décembre 2022.

\bibitem[Fabas, 2023]{fabas2023}
FABAS Pierre-Jean. « L'histoire des réformes de la retraite depuis 1982 ». \textit{LégiSocial}, 24 mars 2023.

\bibitem[La Retraite en clair, s.d.]{retraiteenclair}
LA RETRAITE EN CLAIR. « Les grandes dates de la retraite en France ». Consulté le 3 novembre 2026.

\bibitem[Saly, 2020]{saly2020}
SALY Pierre. « Capitalisation et répartition dans l'histoire du système français des retraites avant 1945 ». \textit{Revue d'histoire de la protection sociale}, n° 13, 2020/1, p. 23-35.

\bibitem[Universalis, 1995]{universalis1995}
ENCYCLOPÆDIA UNIVERSALIS. « 3-31 décembre 1995 : concessions du gouvernement et fin des grèves sur fond d'insatisfaction des partenaires sociaux ». Consulté le 3 novembre 2026.

\bibitem[Universalis, 2010]{universalis2010}
ENCYCLOPÆDIA UNIVERSALIS. « 2-28 octobre 2010 : adoption de la réforme des retraites ». Consulté le 3 novembre 2026.

\bibitem[Valat, 2022]{valat2022}
VALAT Bruno. « Les retraites et la création de la Sécurité sociale en 1945 : révolution ou restauration ? ». Intervention devant le Conseil d'orientation des retraites, 15 décembre 2022.

% \bibitem{auteur_livre}
% NOM Prénom. \textit{Titre du livre en italique}. Lieu d'édition : Nom de l'éditeur, Année.

% \bibitem{auteur_article}
% NOM Prénom. « Titre de l'article entre guillemets ». \textit{Titre de la revue en italique}, volume X, numéro Y, saison Année, p. 123-456.
\end{thebibliography}
